\documentclass{ufctex}

\ufcsetup{
  tipo = tese,
  impressao = anverso,
  ies = {Universidade Federal do Ceará},
  centro = {Centro de Ciências},
  programa-doutorado = {Programa de Pós-Graduação em Ciência da Computação},
  nome-doutorado = {Doutorado em Ciência da Computação},
  titulo-doutor = {Ciência da Computação},
  area-doutorado = {Computação Gráfica},
  autor = {Nome Sobrenome},
  titulo = {Título de Teste da Tese},
  local = {Fortaleza},
  ano = {2026},
  data-aprovacao = {18 de agosto de 2026},
  orientador = {Prof. Dr. Orientador Exemplo},
  orientador-ies = {Universidade Federal do Ceará},
  orientador-unidade = {Centro de Ciências},
  banca-2 = {Prof. Dr. Membro Externo},
  banca-2-ies = {Instituição Externa},
  autor-epigrafe = {Autor de exemplo},
  ficha-catalografica = nao,
  brasao = sim
}

\begin{document}
\pretextual

\imprimircapa
\imprimirfolhaderosto
\imprimirerrata{tests/fixtures/pretextuais/errata.tex}
\imprimirfolhadeaprovacao
\imprimirdedicatoria{tests/fixtures/pretextuais/dedicatoria.tex}
\imprimiragradecimentos{tests/fixtures/pretextuais/agradecimentos.tex}
\imprimirepigrafe[longa]{tests/fixtures/pretextuais/epigrafe.tex}
\imprimirresumo{tests/fixtures/pretextuais/resumo.tex}
\imprimirabstract{tests/fixtures/pretextuais/abstract.tex}
\imprimirlistadefiguras
\imprimirlistadetabelas
\imprimirlistadeabreviaturasesiglas{tests/fixtures/pretextuais/abreviaturas.tex}
\imprimirlistadesimbolos{tests/fixtures/pretextuais/simbolos.tex}
\imprimirsumario

\textual
\section{Introdução}
Texto de teste da seção textual.

\begin{figure}[h]
  \centering
  \rule{4cm}{2cm}
  \caption{Figura de teste}
\end{figure}

\begin{table}[h]
  \centering
  \caption{Tabela de teste}
  \begin{tabular}{cc}
    A & B \\
    1 & 2
  \end{tabular}
\end{table}

\section{Conclusão}
Texto final da fixture.

\end{document}
